\documentclass{article}
\usepackage[T1]{fontenc}
\usepackage{lmodern}
\usepackage{graphicx}
\usepackage{amsmath,amssymb}
\usepackage[a4paper,margin=2cm]{geometry}
\usepackage[absolute,overlay]{textpos}
\setlength{\TPHorizModule}{1mm}
\setlength{\TPVertModule}{1mm}
\usepackage[indonesian]{babel}
\usepackage{hyperref}
\setlength{\emergencystretch}{2em}
% unit: Halaman 29

\begin{document}

% === Halaman 29 (python/native) ===

29

Contoh 9.8. Tinjau kembali plainteks dari Contoh 9.6: 


10100010001110101001 

Bagi plainteks menjadi blok-blok yang berukuran 4 bit: 


1010  0010  0011  1010  1001 


atau  dalam notasi HEX adalah A23A9.  

Misalkan kunci (K) yang digunakan adalah (panjangnya juga 4 bit) 



1011 


atau dalam notasi HEX adalah B. Sedangkan IV yang digunakan seluruhnya bit 0 (Jadi, C0 = 

0000) 

EK(P) = (P  K) << 1

Fungsi enkripsi E yang digunakan sama seperti sebelumnya: XOR-kan blok plainteks Pi dengan 

K, kemudian geser secara wrapping bit-bit dari Pi  K satu posisi ke kiri.

\end{document}
